\documentclass{workbook}

\newcommand{\copyrightdate}{2015--2024}

\usepackage{ifxetex}
\usepackage[utf8]{inputenc}
\usepackage{hyperref}
\usepackage{hyperxmp} % Embed meta data into the PDF
\hypersetup{%
	hidelinks=true,
	linkcolor = {0 0 1},
	% Metadata to be embedded by hyperxmp
	pdftitle={Ordinary Differential Equations (\jobname)},
	pdfauthor={Jason Siefken},
	pdfauthortitle={Author},
	pdfcopyright={Copyright (C) \copyrightdate, Jason Siefken},
	pdfsubject={Ordinary Differential Equations textbook/workbook},
	pdfkeywords={ordinary differential equations, ODEs, modeling, vectors, mathematics, textbook},
	pdfurl={https://github.com/siefkenj/IBLODEs/},
	pdflicenseurl={https://creativecommons.org/licenses/by-sa/4.0/},
}

%%%
% import all needed packages and macros
%%%
\usepackage[yyyymmdd]{datetime}
\input{common/preamble.tex}

\usepackage{breqn}

%%%
% Set up the footers to have the correct copyright notices
%%%

\fancypagestyle{siefken}{%
	\rfoot{\footnotesize\it \copyright\,Jason Siefken, \copyrightdate \ \makebox(30,5){\includegraphics[height=1.2em]{by-sa.pdf}}}
	\lfoot{}
	\renewcommand{\headrulewidth}{0pt}
}
\fancypagestyle{iola}{%
	\rfoot{\footnotesize\it \copyright\,IOLA Team \url{iola.math.vt.edu} \ \makebox(30,5){\includegraphics[height=2.2em]{images/iolalogo.png}}}
	\lfoot{}
	\renewcommand{\headrulewidth}{0pt}
}

\DeclareDocumentEnvironment{iola}{o}{%
	\newpage
	\pagestyle{iola}
}{%
	\newpage
}


%%
% Allow hiding of environments
%%
\usepackage{environ}% http://ctan.org/pkg/environ
\makeatletter
\newcommand{\voidenvironment}[1]{%
  \expandafter\providecommand\csname env@#1@save@env\endcsname{}%
  \expandafter\providecommand\csname env@#1@process\endcsname{}%
  \@ifundefined{#1}{}{\RenewEnviron{#1}{}}%
}
\makeatother
% allow pagebreaks that only display in `standard` mode
\newcommand{\displayonlynewpage}{\begin{displayonly}\newpage\end{displayonly}}
% allow pagebreaks that only display in `book` mode
\newcommand{\bookonlynewpage}{\begin{bookonly}\newpage\end{bookonly}}


%
% Set up the three render modes: standard, instructor, and solutions.
% These render with varying amounts of extra data (like solutions and notes)
%
\newtoggle{instructor}
\newtoggle{standard}
\newtoggle{solutions}
\newtoggle{book}
\newtoggle{slides}
\newcommand{\setinstructor}{
	\toggletrue{instructor}
	\togglefalse{standard}
	\togglefalse{solutions}
	\togglefalse{book}
	\togglefalse{slides}
}
\newcommand{\setstandard}{
	\togglefalse{instructor}
	\toggletrue{standard}
	\togglefalse{solutions}
	\togglefalse{book}
	\togglefalse{slides}
}
\newcommand{\setsolutions}{
	\togglefalse{instructor}
	\togglefalse{standard}
	\toggletrue{solutions}
	\togglefalse{book}
	\togglefalse{slides}
}
\newcommand{\setbook}{
	\togglefalse{instructor}
	\togglefalse{standard}
	\togglefalse{solutions}
	\toggletrue{book}
	\togglefalse{slides}
}
\newcommand{\setslides}{
	\togglefalse{instructor}
	\togglefalse{standard}
	\togglefalse{solutions}
	\togglefalse{book}
	\toggletrue{slides}
}


%
% Infer the document level from the \jobname
%
\usepackage{xstring}
\IfSubStr{\jobname}{\detokenize{book}}{\setbook}{
	\IfSubStr{\jobname}{\detokenize{solutions}}{\setsolutions}{
		\IfSubStr{\jobname}{\detokenize{instructor}}{\setinstructor}{
			\IfSubStr{\jobname}{\detokenize{slides}}{\setslides}{
				\setstandard
			}
		}
	}
}


\setbookoptions{
	twosided = false,
	inline solutions = false,
}


\NewColoredEnvironment{
	name = lesson,
	display name = Lesson,
	banner color = Plum,
	title color = Plum,
	banner on left = true,
	open right = false,
}
\NewColoredEnvironment{
	name = module,
	display name = Module,
	banner color = Turquoise,
	title color = Cerulean,
	definition color = Cerulean,
	theorem color = myorange,
}
\NewColoredEnvironment{
	name = appendix,
	display name = Appendix,
	banner color = LimeGreen,
	title color = LimeGreen!70!Green!80!black,
	definition color = Cerulean,
	theorem color = myorange,
}
\NewColoredEnvironment{
	name = indices,
	display name = Indices,
	banner color = Green,
	title color = Green,
}
\NewColoredEnvironment{
	name = tutorial,
	display name = Tutorial,
	banner color = Peach,
	title color = Peach!80!black,
	emphbox color = Peach,
	% We will print tutorial worksheets back-to-back to save space
	open right = false,
}




\loadgeometry{default}

%
% Hide the non-problem environments
%
\newcommand{\coversubtitle}{} % we override the subtitle in each mode, so make sure the command exists to override.
\iftoggle{instructor}{
	\voidenvironment{module}
	\voidenvironment{appendix}
	\voidenvironment{bookonly}
	\voidenvironment{displayonly}
	\renewcommand{\coversubtitle}{Instructor Guide}
}{}
\iftoggle{solutions}{
	\voidenvironment{module}
	\voidenvironment{appendix}
	\voidenvironment{bookonly}
	\voidenvironment{displayonly}
	\voidenvironment{lesson}
	\voidenvironment{notes}
	\renewcommand{\coversubtitle}{Solutions}
}{}
\iftoggle{standard}{
	\voidenvironment{module}
	\voidenvironment{appendix}
	\voidenvironment{bookonly}
	\voidenvironment{solution}
	\voidenvironment{annotation}
	\voidenvironment{lesson}
	\renewcommand{\coversubtitle}{MAT223 Notes}
	\loadgeometry{default}
}{}
\iftoggle{book}{
	\voidenvironment{displayonly}
	\voidenvironment{solution}
	\voidenvironment{annotation}
	\voidenvironment{lesson}
	\renewcommand{\coversubtitle}{{\hspace{-5pt}\begin{tabular}{l}MAT223 Workbook\\\small\today{} Edition\end{tabular}}}
	\setbookoptions{
		twosided = true,
		inline solutions = false,
	}
	\loadgeometry{book}
}{}
\iftoggle{slides}{
	\voidenvironment{module}
	\voidenvironment{appendix}
	\voidenvironment{bookonly}
	\voidenvironment{solution}
	\voidenvironment{annotation}
	\voidenvironment{lesson}
	\renewcommand{\coversubtitle}{MAT223 Slides}
	\loadgeometry{slides}
	\initSlides
}{}
%\voidenvironment{solution}
%\voidenvironment{annotation}
%\voidenvironment{lesson}
%%\voidenvironment{notes}
%%\voidenvironment{displayonly}

% Allow an index to be created
\makeindex[title=Index of Terms, columns=3]
\makeindex[name=definitions, title=Index of Definitions, columns=3]
\makeindex[name=symbols, title=Index of Symbols, columns=3]

\indexsetup{
	level=\Heading,
	noclearpage
}

\begin{document}
%%
%% Import definitions from definition.tex; all definitions can be restated multiple times
%%

\input{common/definitions.tex}

%%
%% End Definitions
%%




% Needed to get different PDF bookmarks from the TOC entries
\hypersetup{bookmarksdepth=3}

\pagestyle{empty}

\input{common/cover.tex}

\newpage

\begin{bookonly}
	\clearpage
	\hbox{}
	\newpage
	\input{modules/preface.tex}
	\section*{Contributors}
	\input{common/contributors.tex}
	\section*{Dedication}
	\begin{center}
		This book is dedicated to
		\href{https://www.gazettetimes.com/news/local/obituaries/dr-robert-main-burton/article_9c087f07-c005-515a-bb3f-2c9c6a6b7332.html}{\color{blue}Dr.~Bob Burton}---friend and mentor.

		\emph{\large ``Sometimes you have to walk the mystical path with practical feet.''}
	\end{center}
	\newpage
	\mbox{}
	{
		\pagestyle{empty}
		\setcounter{tocdepth}{1}
		\tableofcontents
		\thispagestyle{empty}
	}
	\newpage
	\mbox{}
	\newpage
\end{bookonly}

\setcounter{page}{1}
\pagestyle{siefken}


\addcontentsline{toc}{chapter}{Lessons}

\begin{module}\label{module1}
	\Title{Sets, Vectors \& Notation}

	In this module you will learn
	\begin{itemize}
		\item The basics of sets and set-builder notation.
		\item The definition of vectors and how they relate to points.
		\item Column vector notation and how to represent vectors in drawings.
		\item How to compute linear combinations of vectors and use systems of linear
			equations to answer questions about linear combinations of vectors.
	\end{itemize}

	\Heading{Modeling}

Suppose you are observing some \emph{green} ants walking on the sidewalk.
In the first minute you record 10 ants. In the second minute you
record 20. In the third minute, you record 40 ants. This continues until
there are too many ants for you to count.

\begin{center}
	\begin{tabular}{c|c}
		Minute & \#Green Ants\\
		\hline
		1 & 10\\
		2 & 20\\
		3 & 40\\
		4 & 80\\
		$\vdots$ & $\vdots$
	\end{tabular}
\end{center}

Since you lost count of the ants, you decide to use mathematics to try and figure out
how many ants walked by on minutes $5$, $6$, \ldots. You notice the pattern that
\[
	\text{Green ants per minute $n$} = 2^{n-1}\cdot 10.
\]
Stupendous! Mathematics now predicts there were $160$ ants during minute $5$. But something
else catches your eye. Across the sidewalk are \emph{brown} ants. You count these
ants every minute.

\begin{center}
	\begin{tabular}{c|c}
		Minute & \#Brown Ants\\
		\hline
		1 & 3\\
		2 & 6\\
		3 & 12\\
		4 & 24\\
		$\vdots$ & $\vdots$
	\end{tabular}
\end{center}

The pattern is slightly different. This time, 
\[
	\text{Green ants per minute $n$} = 2^{n-1}\cdot 3.
\]

Your friend, who was watching you the whole time, looks confused. ``Why come up with two complicated equations
when you can describe both types of ant at once?'' they declare.

\begin{center}
	\begin{tabular}{c}
		$\text{\#Ants at minute $n$}\ =\ 2\cdot(\text{\#Ants at minute $n-1$})$\\
		$\text{\#Green ants at minute 1}=10$\\
		$\text{\#Brown ants at minute 1}=3$\\
	\end{tabular}
\end{center}

Your friend has a point. Their model is elegant, but \emph{your} model can predict how many ants pass by at minute $3.222$!
Though, your friend would probably complain that $46.654$ is not a number of ants\ldots.

\medskip

You and your friend have just come up with two different \emph{mathematical models} for the number of ants
that walk across the sidewalk. They happen to make similar predictions for each minute and each have their
strengths and weakenesses. In this course, we will be focused on a particular type of mathematical model---one
that uses \emph{differential equations} at its core.

\Heading{Types of Models}

\begin{definition}[Mathematical Model]
A \emph{mathematical model}\index{model} is a description of the world
	\begin{enumerate}
		\item created in the service of answering a question, and
		\item where the complexity of the world has been abstracted away to numbers, quantities, and their relationships\footnote{ Other
mathematical objects are also allowed.}.
	\end{enumerate}
\end{definition}

In the previous situation, the \emph{question} you were trying to answer was ``how many ants are there at a given minute?''.
We sidestepped difficult issues like, ``Is an ant that is missing three legs still an ant?'' by using the common-sense
convention that ``the number of ants is a whole number and one colored blob that moves under its own power corresponds to one ant''; thus,
we could use single numbers to represent our quantity of interest (the ants).

You and your friend already came up with two types of models.
\begin{itemize}
	\item An \textbf{analytic} model based on known functions.
	\item A \textbf{recursive} model where subsequent terms are based on previous terms and initial conditions.
\end{itemize}

If we define $A(n)$ to be the number of ants corssing the sidewalk at minute $n$, the \emph{analytic} model presented for green ants is
\[
	A(n)=2^{n-1}\cdot 10
\]
and the \emph{recursive} model presented is
\begin{align*}
	A(1) &= 10\\
	A(n) &= 2\cdot A(n-1).
\end{align*}

Each type of model has pros and cons. For example, the analytic model allows you to calculate the number of ants at any minute
with few button presses on a calculator, whereas the recursive model is more difficult to calculate but
makes it clear that the number of ants is doubling every minute.

Often times recursive models are easier to write down than analytic models\footnote{ In fact, in many real-world situations, an analytic model
doesn't exist}, but they maybe harder to analyze. A third type of model has similarities to both analytic and recursive models, and
brings the power of calculus to modeling.

\begin{itemize}
	\item A \textbf{differential-equations} model is a model based on a relationship between a function's derivative(s), its values, and initial conditions.
\end{itemize}

\emph{Differential-equations} models are useful because derivatives correspond to rates of change---and things in the world are always changing.
Let's try to come up with a differential equations model for the ants.

We'd like an equation relating $A(n)$, the number of ants at minute $n$, to $A'(n)$, the \emph{instantaneous rate of change} of the number of ants at minute $n$.
Making a table, we see

\begin{center}
	\begin{tabular}{c|c|c}
		Minute & \#Brown Ants & Change (from prev. minute)\\
		\hline
		1 & 3 & ?\\
		2 & 6 & 3\\
		3 & 12 & 6\\
		4 & 24 & 12\\
	\end{tabular}
\end{center}

or

\begin{center}
	\begin{tabular}{c|c|c}
		Minute & \#Brown Ants & Change (from next minute)\\
		\hline
		1 & 3 & 3\\
		2 & 6 & 6\\
		3 & 12 & 12\\
		4 & 24 & ?\\
	\end{tabular}
\end{center}

depending on whether we record the change from the previous minute or up to the subsequent minute. Neither table gives the \emph{instantaneous} rate of
change, but in both tables, the change is proportional to the number of ants. So, we can set up a model
\[
	A'(n) = kA(n)
\]
where $k$ is a constant of proportionality that we will try to determine later. We've just written down a \emph{differential equation} with an undetermined parameter, $k$.

\begin{definition}[Differential Equation]
	A \emph{differential equation}\index{differential equation} is an equation relating a function to one or more of its derivatives.
\end{definition}

We'd like to figure out what $k$ is. One way to do so is to solve the differential equation and find the values of $k$ so that our model
correctly predicts the data. This is called \emph{fitting} the model to data.

\begin{definition}[Fitting a Model]
	Given a model $M$ with parameters $k_1$, $k_2$, $\ldots$ and data $D$, \emph{fitting the model $M$ to the data $D$}
	is the process of finding values for the parameters $k_1$, $k_2$, $\ldots$ so that $M$ most accurately predicts the data $D$.
\end{definition}

Note that, in general, fitting a model to data doesn't necessarily produce \emph{unique} values for the unknown parameters, and a fitted model
(especially when the data comes from real-world observations) usually doesn't reproduce the data exactly. However, in the case of these ants, we
just might get lucky.

\Heading{Solving Differential Equations}

In general, \emph{there is no algorithm for solving differential equations}. Fortunately, it is easy to check whether any particular
function is a solution to a differential equation, since there \emph{is} an algorithm to differentiate functions\footnote{ More specifically, there is an algorithm
to differentiate the \emph{elementary} functions, those functions formed by compositions, sums, products, and quotients of polynomials, trig, exponentials, and logs.}.
Because of this, \emph{guess and check} will be our primary method for solving differential equations.

\begin{example}
Use educated guessing to solve $A'(n)=kA(n)$.

Since $A'\approx A$, we might start with a function that is equal to its own derivative. There is a famous one: $e^n$. Testing, we see
\[
	\frac{\d}{\d n}e^n=e^n=ke^n
\]
if $k=1$, but it doesn't work for other $k$'s. Trying $e^{kn}$ instead yields
\[
	\frac{\d}{\d n}e^{kn}=ke^{kn}
\]
which holds for all $k$. Thus $A(n)=e^{kn}$ is \emph{a} solution to $A'(n)=kA(n)$. However, there are other solutions, because
\[
	\frac{\d}{\d n}Ce^{kn}=C\left(ke^{kn}\right)=k\left(Ce^{kn}\right),
\]
and so for every $C$, the function $A(n)=Ce^{kn}$ is a solution to $A'(n)=kA(n)$.
\end{example}

By guessing-and-checking, we have found an infinite number of solutions to $A'(n)=kA(n)$.  It's now time to fit our
model to the data.

\begin{example}
	Find values of $C$ and $k$ so that $A(n)=Ce^{kn}$ best models brown ants.

	Taking two rows from our brown ants table, we see
	\begin{align*}
		A(1) &= Ce^{k} = 3\\
		A(2) & = Ce^{2k}=6.
	\end{align*}

	Since $e^k$ can never be zero, from the first equation we get $C=3/e^k$. Combining with the second equation we find
	\[
		Ce^{2k}=\frac{3}{e^k}e^{2k}=3e^k=6
	\]
	and so $e^k=2$. In other words $k=\ln 2$. Plugging this back in, we find $C=3/2$. Thus our fitted model is
	\[
		A(n)=\tfrac{3}{2}e^{n\ln 2}.
	\]
\end{example}

Upon inspection, we can see that $\tfrac{3}{2}e^{n\ln 2} = 3\cdot 2^{n-1}$, which is the analytic model
that was first guessed for brown ants.



%One of the beauties of the guess-and-check method is that it \emph{cannot} produce an incorrect answer,
%since you're always checking (though it may fail to produce answers, any answer that it \emph{does} produce
%will be correct). Thus, even if you did some mathematically-dubious steps to produce a guess, if it passes the checks,
%the final answer is valid.
%
%\begin{example}
%Use Leibniz notation, integral calculus, and guess-and-check to come up with solutions to $A'(n)=kA(n)$.
%
%In Leibniz notation, $A'(n)$ is written as $\displaystyle\frac{\d A}{\d n}$, and we may suppress the function variable and write $A$ in place of $A(n)$. Thus, our equation becomes
%\[
%	\frac{\d A}{\d n}=kA.
%\]
%Treating $\d A/\d n$ as an actual ratio, we can bring all $A$'s to the left and $n$'s to the right, giving
%\[
%	\frac{1}{A}\d A = k\d n.
%\]
%Integrating both sides gives
%\[
%	\int \frac{1}{A}\d A = \ln A + C_1 \qquad =\qquad  \int k\d n = kn + C_2,
%\]
%and so
%\[
%	\ln A + C_1 = kn + C_2 \qquad \iff \qquad \ln A = kn + (C_2-C_1) = kn + C_3.
%\]
%Solving for $A$ by exponentiating both sides gives
%\[
%	A = e^{\ln A} = e^{kn + C_3} = C_4 e^{kn}.
%\]
%Much of what we've done in the preceding calculation \emph{was not rigorous mathematics}, but that's okay, because
%it leads us to the guess: $A(n)=Ce^{kn}$. Finally, we can test our guess and verify
%\[
%	\frac{\d}{\d n}Ce^{kn}=C\left(ke^{kn}\right)=k\left(Ce^{kn}\right),
%\]
%and so $A(n)=Ce^{kn}$ is a solution for every $C$.\footnote{ Notice that in our ``guess work'', $C_4$ was actually restricted to be a positive
%number, but our solution works with $C$ being positive or negative.}
%
%\end{example}




	\input{modules/module1-exercises.tex}
\end{module}

\begin{lesson}
	\Title{Linear Combinations}

	\Heading{Textbook}
	Module 1

	\Heading{Objectives}
	\begin{itemize}
		\item Internalize vectors as geometric objects representing displacements.

		\item Use column vector notation to write vectors.

		\item Relate points and vectors and be able to interpret a point as
			a vector and a vector as a point.

		\item Solve simple equations involving vectors.
	\end{itemize}

	\Heading{Motivation} Students have differing levels of experience with vectors.
	We want to establish a common notation for vectors and use vector notation
	along with algebra to solve simple questions. E.g., ``How can I get to location
	$A$ given that I can only walk parallel to the lines $y=4x$ and $y=-x$?''


	\begin{annotation}
		\begin{notes}
			\begin{itemize}
			\item
			We will use the language \emph{component of $\vec v$ in
			the direction $\vec u$} in the future and it will be a \emph{vector}.
			For this reason, try to refer to the entries of a column
			vector as \emph{coordinates} or \emph{entries} instead of components.

			\item
			Though we will almost exclusively use
			column vector notation in this course, students should be able to parse
			questions phrased in terms of row vectors.
			\end{itemize}
		\end{notes}
	\end{annotation}We will use column vector notation and the idea of equating
	coordinates in order to solve problems.

\end{lesson}

\begin{iola}
\begin{slide}
\section*{The Magic Carpet Ride}
\addcontentsline{toc}{subsection}{The Magic Carpet Ride}

\question
\begin{annotation}
	\begin{goals}
		\Goal{Hands-on experience with vectors as displacements.}
		\begin{itemize}
			\item Internalize vectors as geometric objects representing
				displacements.

			\item Use column vector notation to write vectors.

			\item Use pre-existing knowledge of algebra to answer vector
				questions.
		\end{itemize}
	\end{goals}
	\begin{notes}

		\begin{itemize}
			\item There are many ways to solve this problem.
				Some students
				might start with equations. After they use their
				equations to solve the problem, make them draw a picture
				and come up with a graphical solution.

			\item When the students start coming up with vector equations,
				give them the vocabulary of \emph{linear
				combinations}
				and \emph{column vector notation}.
		\end{itemize}
	\end{notes}
\end{annotation}
You are a young adventurer. Having spent most of your time in the mythical city of Oronto,
	you decide to leave home for the first time. Your parents
want to help you on your journey, so just before your departure, they give you two
gifts. Specifically, they give you two forms of transportation: a hover board and
a magic carpet. Your parents inform you that both the hover board and the magic carpet
have restrictions in how they operate:

\begin{minipage}{\linewidth}
	\vspace{.5cm}
	\begin{wrapfigure}{l}{1in}
	\vspace{-.8cm}
	\includegraphics[width=1in]{images/HoverBoard-small.png}
	\end{wrapfigure}

	We denote the restriction on the hover board's movement by the vector
	$\mat{3 \\1}$. By this we mean that if
	the hover board traveled ``forward'' for one hour, it would move along a
	``diagonal'' path that would result in a displacement of 3 km East and
	1 km North of its starting location.
\end{minipage}

\begin{minipage}{\linewidth}
	\vspace{.5cm}
	\begin{wrapfigure}{l}{1in}
	\vspace{-.8cm}
	\includegraphics[width=1in]{images/MagicCarpet2-small.png}
	\end{wrapfigure}

	We denote the restriction on the magic carpet's movement by the vector
	$\mat{1 \\2 }$. By this we mean that if the
	magic carpet traveled ``forward'' for one hour, it would move along a
	``diagonal'' path that would result in a displacement of 1 km East and
	2 km North of its starting location.
\end{minipage}

\lfoot{\footnotesize Drawings by \url{@DavidsonJohnR} (twitter)}

\vspace{10mm}
\end{slide}

\begin{slide}
% Scenario Section
\textbf{Scenario One: The Maiden Voyage}

Your Uncle Cramer suggests that your first adventure should be to go visit
the wise man, Old Man Gauss. Uncle Cramer tells you that Old Man Gauss
lives in a cabin that is 107 km East and 64 km North of your home.

\vspace{5mm}

\textbf{Task:}
\par
Investigate whether or not you can use the hover board and the magic
carpet to get to Gauss's cabin. If so, how? If it is not possible to
get to the cabin with these modes of transportation, why is that the case?

%\vspace{5mm}
% As a group, state and explain your answer(s) on the group whiteboard. Use
% the vector notation for each mode of transportation as part of your
% explanation and use a diagram or graphic to help illustrate your
% point(s).

\end{slide}
\end{iola}

\begin{lesson}
	\Title{Linear Combinations}

	\Heading{Textbook}
	Module 1

	\Heading{Objectives}
	\begin{itemize}
		\item Set up and solve vector equations $a\vec v+b\vec u=\vec w$. The solving
			method may be ad hoc.
		\item Use set notation and set operations/relations $\cup$, $\cap$, $\in$, $\subseteq$.
		\item Translate between set-builder notation and words in multiple ways.
	\end{itemize}

	\Heading{Motivation}
	We revisit questions about linear combinations more formally and generate a need for
	algebra. The algebra we do to solve vector equations will become algorithmic when
	we learn row reduction, but at the moment, any method is fine.

	\begin{annotation}
		\begin{notes}
			You will have a mix of MAT135/136 and MAT137 students.
			The MAT137 students will be doing logic and sets in their
			class. The MAT135 students won't. Make sure not to leave them
			behind!
		\end{notes}
	\end{annotation}
	As we talk about more complex objects, we need precise ways to talk about
	groups of vectors. I.e., we need sets and set-builder notation. This preview of set-builder
	notation will take some of the difficulty away when we define span as a set of vectors.

	In this course we will be using formal and precise language. Part of this lesson
	is that there are multiple correct ways (and multiple incorrect ways) to use formal
	language. Gone are the days of ``there's only one right answer and it is 4''!

\end{lesson}

\begin{iola}
\begin{slide}
\section*{The Magic Carpet Ride, Hide and Seek}
\addcontentsline{toc}{subsection}{The Magic Carpet Ride, Hide and Seek}
\question

\begin{annotation}
	\begin{goals}
		\Goal{Address an existential question involving vectors: ``Is it possible
		to find a linear combination that does\ldots?''}

		The goal of this problem is to
		\begin{itemize}
			\item Formalize geometric questions using the language of vectors.
			\item Find both geometric and algebraic arguments to support the same
				conclusion.
			\item Establish what a ``negative multiple'' of a vector should be.
		\end{itemize}
	\end{goals}
	\begin{notes}
		\begin{itemize}
			\item Both \emph{yes} and \emph{no} are valid answers to
				this question depending on whether you are allowed
				to go backwards. Establish that ``negative'' multiples of
				a vector mean traveling backwards along that vector.
			\item This problem can be solved with algebra by finding a formula
				for the coefficients for an arbitrary position or with geometry,
				with arguments eventually hinging on the fact that non-parallel
				lines intersect.
		\end{itemize}
	\end{notes}
\end{annotation}
You are a young adventurer. Having spent most of your time in the mythical city of Oronto,
	you decide to leave home for the first time. Your parents
want to help you on your journey, so just before your departure, they give
you two gifts. Specifically, they give you two forms of transportation:
a hover board and a magic carpet. Your parents inform you that both the
hover board and the magic carpet have restrictions in how they operate:



\begin{minipage}{\linewidth}
	\vspace{.5cm}
	\begin{wrapfigure}{l}{1in}
	\vspace{-.8cm}
	\includegraphics[width=1in]{images/HoverBoard-small.png}
	\end{wrapfigure}

	We denote the restriction on the hover board's movement by the vector
	$\mat{3 \\1}$. By this we mean that if
	the hover board traveled ``forward'' for one hour, it would move along a
	``diagonal'' path that would result in a displacement of 3 km East and
	1 km North of its starting location.
\end{minipage}

\begin{minipage}{\linewidth}
	\vspace{.5cm}
	\begin{wrapfigure}{l}{1in}
	\vspace{-.8cm}
	\includegraphics[width=1in]{images/MagicCarpet2-small.png}
	\end{wrapfigure}

	We denote the restriction on the magic carpet's movement by the vector
	$\mat{1 \\2 }$. By this we mean that if the
	magic carpet traveled ``forward'' for one hour, it would move along a
	``diagonal'' path that would result in a displacement of 1 km East and
	2 km North of its starting location.
	\vspace{1cm}
\end{minipage}

\end{slide}
\begin{slide}

\textbf{Scenario Two: Hide-and-Seek}

Old Man Gauss wants to move to a cabin in a different location. You are
not sure whether Gauss is just trying to test your wits at finding him
or if he actually wants to hide somewhere that you can't visit him.

\vspace{5mm}

\textbf{Are there some locations that he can hide and you cannot reach him
with these two modes of transportation?}

Describe the places that you
can reach using a combination of the hover board and the magic carpet and
those you cannot. Specify these geometrically and algebraically. Include
a symbolic representation using vector notation. Also, include a convincing
argument supporting your answer.

%\vspace{5mm} \par \textbf{Use your
%group's whiteboard as a space to write out our work as your work together
%on this problem.}
\end{slide}
\end{iola}





\begin{slide}[+1]

\section*{Sets and Set Notation}
\vspace{-.5cm}

	\SavedDefinitionRender{Set}

	\begin{definition}
		\raggedright
		Some common sets are
		\begin{itemize}
			\item[] $\N=\Set{\text{natural numbers}} = \Set{\text{non-negative whole numbers}}$.
			\item[] $\Z=\Set{\text{integers}} = \Set{\text{whole numbers, including negatives}}$.
			\item[] $\R=\Set{\text{real numbers}}$.
			\item[] $\R^n=\Set{\text{vectors in $n$-dimensional Euclidean space}}$.
		\end{itemize}
	\end{definition}
\end{slide}
\begin{slide}

	\question
	\begin{annotation}
		\begin{goals}
			\Goal{Practice reading sets and set-builder notation.}

			The goal of this problem is to
			\begin{itemize}
				\item Become familiar with $\in$, $\subseteq$, and $=$ in
					the context of sets.
				\item Distinguish between $\in$ and $\subseteq$.
				\item Use quantifiers with sets.
			\end{itemize}
		\end{goals}

		\begin{notes}
			\begin{itemize}
				\item Most are easy up through (h).
				\item Make students ``fix'' (i) so it
					becomes true.
				\item (j) and (k) are an opportunity to use
					the definition of set equality. Students don't
					realize that $=$'s has a definition.
			\end{itemize}
		\end{notes}
	\end{annotation}
	\begin{parts}
		\item Which of the following statements are true?
		\begin{enumerate}
			\item $3\in\Set{1,2,3}$.
				\begin{solution}[inline]True\end{solution}
			\item $1.5\in\Set{1,2,3}$.
				\begin{solution}[inline]False\end{solution}
			\item $4\in\Set{1,2,3}$.
				\begin{solution}[inline]False\end{solution}
			\item ``b''$\in \Set{ x \given x\text{ is an English letter}}$.
				\begin{solution}[inline]True\end{solution}
			\item ``\`o''$\in \Set{x \given x\text{ is an English letter}}$.
				\begin{solution}[inline]False\end{solution}
			\item $\Set{1,2} \subseteq \Set{1,2,3}$.
				\begin{solution}[inline]True\end{solution}
			\item For some $a\in \Set{1,2,3}$, $a \geq 3$.
				\begin{solution}[inline]True\end{solution}
			\item For any $a\in \Set{1,2,3}$, $a\geq 3$.
				\begin{solution}[inline]False\end{solution}
			\item $1 \subseteq \Set{1,2,3}$.
				\begin{solution}[inline]False\end{solution}
			\item $\Set{1,2,3}=\Set{x\in\R \given 1\leq x\leq 3}$.
				\begin{solution}[inline]False\end{solution}
			\item $\Set{1,2,3}=\Set{x\in\Z \given 1\leq x\leq 3}$.
				\begin{solution}[inline]True\end{solution}
		\end{enumerate}
	\end{parts}
\end{slide}
\begin{slide}
	\bookonlynewpage
	\question
	\begin{annotation}
		\begin{goals}
			\Goal{Practice writing sets using set-builder notation.}

			The goal of this problem is to
			\begin{itemize}
				\item Express English descriptions using math notation.
				\item Recognize there is more than one correct way to
					write formal math.
				\item Preview vector form of a line.
			\end{itemize}
		\end{goals}

		\begin{notes}
			\begin{itemize}
				\item There are multiple correct ways to write
					each of these sets. It's a good opportunity
					to get many correct and incorrect sets up on the
					board for discussing.
				\item Don't worry about the geometry of $B$. That's coming
					in a later problem.
			\end{itemize}
		\end{notes}
	\end{annotation}
		Write the following in set-builder notation
	\begin{parts}
			\item The subset $A\subseteq \R$ of real numbers larger than $\sqrt{2}$.
				\begin{solution}
					$\Set*{x\in\R \given x>\sqrt{2}}$.
				\end{solution}
			\item The subset $B\subseteq \R^2$ of vectors whose first coordinate
			is twice the second.
				\begin{solution}
					$\Set*{\vec v\in\R^2\given \vec v=\mat{a\\b}\text{ with }a=2b}$
					or
					$\Set*{\vec v\in\R^2\given\vec v=\mat{2t\\t}\text{ for some }t\in \R}$\\
					or
					$\Set*{\mat{a\\b}\in\R^2\given a=2b}$.
				\end{solution}
	\end{parts}
\end{slide}
\begin{slide}
	\displayonlynewpage
	\begin{bookonly}\vspace{\stretch{1}}\end{bookonly}
	\SavedDefinitionRender{UnionsIntersections}

	\question
	\begin{annotation}
		\begin{goals}
			\Goal{Apply the definition of $\cup$ and $\cap$.}
		\end{goals}

		\begin{notes}
			\begin{itemize}
				\item It's not important to emphasize that $\cup$ and $\cap$ are binary
			operations but we ask for $X\cup Y\cup Z$ without parenthesis.
			Students won't worry if you don't bring it up.
				\item It won't be clear to them how to write the empty set.
					Some will write $\{\emptyset\}$. Make sure this comes out.
			\end{itemize}
		\end{notes}
	\end{annotation}
	Let $X=\Set{1,2,3}$ and $Y=\Set{2,3,4,5}$ and $Z=\Set{4,5,6}$.  Compute
	\begin{parts}
		\item $X\cup Y$ \begin{solution}[inline]$\Set{1,2,3,4,5}$\end{solution}
		\item $X\cap Y$ \begin{solution}[inline]$\Set{2,3}$\end{solution}
		\item $X\cup Y\cup Z$ \begin{solution}[inline]$\Set{1,2,3,4,5,6}$\end{solution}
		\item $X\cap Y\cap Z$ \begin{solution}[inline]$\emptyset=\Set{}$\end{solution}
	\end{parts}
	\begin{bookonly}\vspace{\stretch{1}}\end{bookonly}
\end{slide}

%\begin{module}
%All Defs:
%\edef\mylist{\pgfkeysvalueof{/definitions/all}}
%\foreach \x in \mylist {
%	\x
%	\SavedDefinitionRender{\x}
%}
%\end{module}









\begin{lesson}
	\Title{Visualizing Sets, Formal Language of Linear Combinations}

	\Heading{Textbook}
	Module 1

	\Heading{Objectives}
	\begin{itemize}
		\item Draw pictures of formally-described subsets of $\R^2$.
		\item Graphically represent $\cup$ and $\cap$ for subsets of $\R^2$.
		\item Graphically represent linear combinations and then come up with
			algebraic arguments to support graphical intuition.
	\end{itemize}

	\Heading{Motivation}

	We want to build a bridge between the formal language of linear combinations
	and set-builder notation and geometric intuition. Where as last time
	the focus was on formal language, this time the focus is on linking geometry
	to formal descriptions.


\end{lesson}

\begin{slide}
	\bookonlynewpage
	\question
	\begin{annotation}
		\begin{goals}
			\Goal{Visualize sets of vectors.}

			The goal of this problem is to
			\begin{itemize}
				\item Apply set-builder notation in the context of vectors.
				\item Distinguish between ``for all'' and ``for some''
					in set builder notation.
				\item Practice unions and intersections.
				\item Practice thinking about set equality.
			\end{itemize}
		\end{goals}

		\begin{notes}
			\begin{itemize}
				\item 1--3 will be easy.
				\item Have a discussion about when you
					should draw vectors as arrows
					vs\mbox{.} as points.
				\item 4 gets at a subtle point that will come up again
					when we define span.
				\item Many will miss 7. Writing a proof for this is
					good practice.
			\end{itemize}
		\end{notes}
	\end{annotation}
	Draw the following subsets of $\R^2$.
	\begin{parts}
		\item $V=\Set*{\vec x\in\R^2 \given \vec x=\mat{0\\t}\text{ for some }t\in\R}$.
		\item $H=\Set*{\vec x\in\R^2 \given \vec x=\mat{t\\0}\text{ for some }t\in\R}$.
		\item $D=\Set*{\vec x\in\R^2 \given \vec x=t\mat{1\\1}\text{ for some }t\in\R}$.
		\begin{solution}
	\begin{tikzpicture}[scale=1.2, >=latex]
    \begin{axis}[scale=1,
		    axis equal image,
		    axis line style={draw=none},
		    tick style={draw=none},
		    yticklabels={,,},
		    xticklabels={,,},
		 xmin=-4.5,
		 xmax=4.5,
		 ymin=-4.5,
		 ymax=4.5,
		 major grid style={dotted, gray},
                 xtick={-10,-9,...,10},
                 ytick={-10,-9,...,10},
                 grid=both,
		 anchor=origin]

	    \draw[Green, very thick] (-5,0) -- (5,0) node[near end, below] {$H$};
	    \draw[cyan!80!black, very thick] (0,-5) -- (0,5) node[near end, right] {$V$};
	    \draw[orange!80!black, very thick] (-5,-5) -- (5,5) node[near end, below right] {$D$};

	    \fill[fill=black] (0,0) circle[radius=2pt] node[below right, xshift=5pt] {\color{black}$\vec 0$};
    \end{axis}
\end{tikzpicture}
		\end{solution}
		\item $N=\Set*{\vec x\in\R^2 \given \vec x=t\mat{1\\1}\text{ for all }t\in\R}$.
				\begin{solution}[inline]
			$N=\Set{}$.
		\end{solution}

		\item $V\cup H$.
			\begin{solution}[inline]
			$V\cup H$ looks like a ``$+$'' going through the origin.
		\end{solution}
		\item $V\cap H$.
			\begin{solution}[inline]
				$V\cap H=\Set{\vec 0}$ is just the origin.
		\end{solution}
		\item Does $V\cup H=\R^2$?
			\begin{solution}
				No. $V\cup H$ does not contain $\mat{1\\1}$ while $\R^2$ does contain
				$\mat{1\\1}$.
			\end{solution}
	\end{parts}

	\begin{bookonly}\begin{center}\hspace{-2cm}\triplegrid\hspace{-2cm}\triplegrid\end{center}\end{bookonly}
	\bookonlynewpage
	\displayonlynewpage
\section*{Vector Combinations}
	\vspace{-1em}

\end{slide}
\begin{slide}

	\SavedDefinitionRender{LinearCombination}

	\question
	\label{ProbSkewBasis}
	\begin{annotation}
		\begin{goals}
			\Goal{Practice linear combinations.}

			The goal of this problem is to
			\begin{itemize}
				\item Practice using the formal term \emph{linear combination}.
				\item Foreshadow span.
			\end{itemize}
		\end{goals}

		\begin{notes}
			\begin{itemize}
				\item In 2, the question should arise: ``Is $3\vec v_1$
					a linear combination of $\vec v_1$ \emph{and}
					$\vec v_2$?'' Address this.
				\item Refer to the magic carpet ride for 5. You don't
					need to do a full proof.
			\end{itemize}
		\end{notes}
	\end{annotation}
	Let $\vec v_1=\mat{1\\1}$, $\vec v_2=\mat{1\\-1}$, and $\vec w=2\vec v_1+\vec v_2$.
	\begin{parts}
		\item Write $\vec w$ as a column vector. When $\vec w$ is written as a
			linear combination of $\vec v_1$ and $\vec v_2$, what are the
			coefficients of $\vec v_1$ and $\vec v_2$?
			\begin{solution}
				$\vec w=\mat{3\\1}$; the coefficients are $(2,1)$.
			\end{solution}
		\item Is $\mat{3\\3}$ a linear combination of $\vec v_1$ and $\vec v_2$?
			\begin{solution}[inline]
				Yes. $\mat{3\\3}=3\vec v_1+0\vec v_2$.
			\end{solution}

		\item Is $\mat{0\\0}$ a linear combination of $\vec v_1$ and $\vec v_2$?
			\begin{solution}[inline]
				Yes. $\vec 0=0\vec v_1+0\vec v_2$.
			\end{solution}
		\item Is $\mat{4\\0}$ a linear combination of $\vec v_1$ and $\vec v_2$?
			\begin{solution}[inline]
				Yes. $\mat{4\\0}=2\vec v_1+2\vec v_2$.
			\end{solution}
		\item Can you find a vector in $\R^2$ that isn't a linear combination of
		$\vec v_1$ and $\vec v_2$?
			\begin{solution}
				No. $\mat{1\\0}=\tfrac{1}{2}\vec v_1+\tfrac{1}{2}\vec v_2$ and
				$\mat{0\\1}=\tfrac{1}{2}\vec v_1-\tfrac{1}{2}\vec v_2$.
				Therefore
				\[
					\mat{a\\b}
					= a\mat{1\\0}+b\mat{0\\1}
					= a(\tfrac{1}{2}\vec v_1+\tfrac{1}{2}\vec v_2)
						+b(\tfrac{1}{2}\vec v_1-\tfrac{1}{2}\vec v_2)
					=(\tfrac{a+b}{2})\vec v_1+(\tfrac{a-b}{2})\vec v_2.
				\]
				Therefore any vector in $\R^2$ can be written as linear combinations
				of $\vec v_1$ and $\vec v_2$.
			\end{solution}
		\item Can you find a vector in $\R^2$ that isn't a linear combination of
			$\vec v_1$?
			\begin{solution}
				Yes. All linear combinations of $\vec v_1$ have equal $x$ and
				$y$ coordinates, therefore $\vec w=\mat{2\\1}$ is not a linear
				combination of $\vec v_1$.
			\end{solution}
	\end{parts}
\end{slide}
\begin{slide}

	\bookonlynewpage
	\question
	\begin{annotation}
		\begin{goals}
			\Goal{Practice formal writing.}
		\end{goals}

		\begin{notes}
			\begin{itemize}
				\item Make everyone \emph{write}. They will think
					they can do it, but they will find it hard if
					they try.
			\end{itemize}
		\end{notes}
	\end{annotation}
	Recall the \emph{Magic Carpet Ride} task where the hover board could
	travel in the direction $\vec h=\mat{3\\1}$ and the magic carpet could
	move in the direction $\vec m=\mat{1\\2}$.
	\begin{parts}
		\item Rephrase the sentence \emph{``Gauss can be reached using just the
			magic carpet and the hover board''} using formal mathematical
			language.
			\begin{solution}
				Gauss's location can be written as a linear combination of
				$\vec m$ and $\vec h$.
			\end{solution}
		\item Rephrase the sentence \emph{``There is nowhere Gauss can hide
			where he is inaccessible by magic carpet and hover board''} using
			formal mathematical language.
			\begin{solution}
				Every vector in $\R^2$ can be written as a linear combination
				of $\vec m$ and	$\vec h$.
			\end{solution}
		\item Rephrase the sentence \emph{``$\R^2$ is the set of all linear
			combinations of $\vec h$ and $\vec m$''} using formal mathematical
			language.
			\begin{solution}
				$\R^2=\Set{\vec v\given \vec v=t\vec m+s\vec h\text{ for some }t,s\in \R}$.
			\end{solution}
	\end{parts}
\end{slide}


\begin{appendix}\label{APPSLEI}
	\Title{Systems of Linear Equations I}

	In this appendix you will learn
	\begin{itemize}
		\item What a system of linear equations is.
		\item What the solution set to a system of equations is, and what it means for a system of equations
			to be consistent or inconsistent.
		\item How augmented matrices can be used to solve systems of linear equations.
		\item How to apply row reduction to find a unique solution to a system of linear
			equations and to determine if a system of linear equations is consistent or inconsistent.
	\end{itemize}

	\input{modules/appendix1.tex}
	\input{modules/appendix1-exercises.tex}
\end{appendix}


\begin{appendix}
	\PrintExerciseSolutions
\end{appendix}
	\begin{indices}*
		\Title{Indices}

		\printindex[symbols]

		\bigskip
		\printindex

		\bigskip
		\printindex[definitions]
	\end{indices}
\end{bookonly}

\end{document}
